MODALITÉS DE PARTAGE ET MODIFICATION DU CAHIER DE CALCUL

- La version PDF du cahier de calcul est librement diffusable par les professeurs auprès de leurs élèves.
- Le PDF ou le document imprimé mis à disposition par le professeur doit faire figurer le QR-code figurant page ii du PDF (ce QR-code permet d'accéder au site web du cahier d'entraînement).
- Les participants au projet et les professeurs peuvent le rendre disponible sur leur site web.
- Le PDF et les sources LaTeX seront mis à disposition des membres de l'UPS via la BEDOC de l'UPS.
- Les sources LaTeX ne sont pas diffusables hors BEDOC.
- Les sources sont mises à disposition pour que les professeurs puissent adapter le cahier à leur classe, par exemple en enlevant un exercice, etc. En revanche, les sources ne sont pas mises à disposition pour servir de base à un autre projet.
- Le contenu du cahier de calcul est modifiable et réutilisable pour des usages spécifiques à l'exclusion de tout usage commercial.
- Les noms des participants au projet doivent être indiqués :

Coordination
Colas Bardavid

Participants
Vincent Bayle
Romain Basson
Olivier Bertrand
Ménard Bourgade
Julien Bureaux
Alain Camanes
Mathieu Charlot
Mathilde Colin de Verdière
Keven Commault
Miguel Concy
Rémy Eupherte
Hélène Gros
Audrey Hechner
Florian Hechner
Marie Hézard
Nicolas Laillet
Valérie Le Blanc
Thierry Limoges
Quang-Thai Ngo
Xavier Pellegrin
Fabien Pellegrini
Jean-Louis Pourtier
Valérie Robert
Jean-Pierre Técourt
Guillaume Tomasini
Marc Tenti